% Imporditud: tools/import_eksperimendid.py
% Allikas: Eesti füüsikaolümpiaadi eksperimendi ülesannete
%   kogu (Rael Kalda, 2023), ülesanne Ü158, lahendus L158.
\setAuthor{Konstantin Dukatš}
\setRound{lõppvoor}
\setYear{2021}
\setNumber{G E1}
\setDifficulty{4}
\setTopic{dünaamika}
\setVahendid{üks A4 paberileht, 1L täis mahlakarp, joonlaud}
\setPunktid{10}
\setAllikas{Eksperimendi kogu Ü158, lahendus lk 119}
\setLahendusseis{toores}

\prob{Mahlakarp}
Leidke mahlakarbi ja paberilehe vaheline hõõrdetegur. Keelatud on kasutada mistahes kaldpinna põhimõtet.

\hint

\solu
Paneme paberilehe lauale ning selle peale asetame mahlakarbi (alusega a $\times$ a). Kui me lükkame joonlauaga horisontaalselt karbi põhja juurest, siis hakkab karp libisema. Kui me aga lükkame karbi ülemist osa, hakkab karp pöörlema. Mingil kõrgusel h on piirjuht, kus karp nii pöörleb kui ka libiseb. Sellel piirjuhul kehtib nii jõudude kui ka jõumomentide tasakaal. Jõudude tasakaalust saame, et joonlauaga lükkamise jõud F = $\mu$mg. Joonlauaga lükkamisel tekkivat jõumomenti tasakaalustab karbile mõjuva gravitatsioonijõu poolt tekitatud jõumoment:

a F h = mg .

Kombineerides neid võrandeid saame, et a

$\mu$= . 2h

Ühe konkreetse AURA mahlakarbi ja paberilehe jaoks olid mõõtmistulemused järgmised: a = 7,0 cm, h = 11,5 cm ning $\mu$ = 0,30. Kuna mahlakarbid ja paberilehed erinevad vähesel määral, siis hõõrdetegur peaks enamasti olema vahemikus 0,28-- 0,42.
\probend
